%----------------------------------------------------------------------------------------
% Tailored CV — PhD position: Genome Stability and Nucleotide Metabolism in Cancer
% Carnie Lab, Klinik Kinderheilkunde III / Hopp Children's Cancer Center (KiTZ), UKHD
% Job-ID V000015650
%----------------------------------------------------------------------------------------

\documentclass[10pt,a4paper,sans]{moderncv}

\usepackage[utf8]{inputenc}
\usepackage[T1]{fontenc}

\moderncvstyle{classic}
\moderncvcolor{blue}

\usepackage[scale=0.78]{geometry}
\usepackage{ragged2e}

%----------------------------------------------------------------------------------------

\firstname{Kundai Farai}
\familyname{Sachikonye}

\address{Auf der Lichtung 19}{82178 Puchheim, Germany}
\mobile{(+49) 1626386344}
\email{kundai.sachikonye@bitspark.com}
\homepage{github.com/fullscreen-triangle}{github.com/fullscreen-triangle}

\quote{Experimental DNA-damage background and computational omics analysis ---
the wet-lab and dry-lab sides of genome stability and metabolism in cancer.}

%----------------------------------------------------------------------------------------

\begin{document}

\makecvtitle

%----------------------------------------------------------------------------------------
%	EDUCATION
%----------------------------------------------------------------------------------------

\section{Education}

\cventry{2019--2021}{PhD candidate, Computational Lipidomics}{Technische Universität München / Bayerisches Forschungsinstitut}{}{}{
  HPC mass-spectrometry pipelines ($>$100\,GB mzML), multi-omics statistical
  modelling, federated analysis across distributed clinical sites. Departed to
  pursue independent computational research.
}
\cventry{2018--2019}{MSc Computational Biology}{Universität Tübingen}{}{}{
  Sequence analysis, variant calling, RNA-seq quantification, pathway analysis;
  bioinformatics pipeline development.
}
\cventry{2014--2017}{MSc Pharmaceutical Biotechnology}{HAW Hamburg}{}{}{
  Molecular biotechnology, drug development, and analytical methods.
}
\cventry{2011--2014}{BSc Biochemistry and Cell Biology}{Jacobs University Bremen}{}{}{
  Thesis: \emph{DNA interstrand crosslinking} --- experimental study of
  crosslink formation and the cellular DNA-damage response; mammalian cell
  culture and fluorescence microscopy.
}

%----------------------------------------------------------------------------------------
%	EXPERIENCE
%----------------------------------------------------------------------------------------

\section{Experience}

\cventry{2021--present}{Independent Computational Biology Researcher \& Developer}{Bitspark GmbH}{}{}{
  Full-time development of validated scientific data pipelines and analysis
  frameworks across genomics, proteomics, metabolomics, and computational
  oncology. All systems publicly documented and deployed.
}

\cventry{2019--2021}{Computational Lipidomics}{TUM / Bayerisches Forschungsinstitut}{Munich}{}{
  MS-based lipidomics at HPC scale: peak detection, ion annotation (LIPID MAPS,
  LipidBlast, HMDB), multi-omics integration, federated processing across
  clinical sites. $>$100\,GB mzML dataset management.
}

\cventry{2017--2018}{Glycomics and Proteomics}{Universitätsklinikum Hamburg-Eppendorf}{Hamburg}{}{
  Automated LC-MS/MS and MALDI-TOF workflows for proteomics and glycomics;
  standardised SOP development for data acquisition and processing.
}

\cventry{2013}{Cancer Biomarker Discovery}{University Medical Center Groningen}{Groningen}{}{
  UPLC-MS/MS shotgun proteomics on tumour samples; data normalisation and
  statistical analysis.
}

%----------------------------------------------------------------------------------------
%	SOFTWARE AND COMPUTATIONAL TOOLS
%----------------------------------------------------------------------------------------

\section{Software \& Computational Tools}

\cvitem{gospel}{
  \textbf{Genomic Variant \& Multi-Omics Analysis.}
  $10^6$ variants/second VCF processing; RNA-seq integration; Bayesian-network
  inference and Gaussian mixture models; pathway analysis (KEGG, Reactome).
  Validated against 1000~Genomes Ph.\ 3, gnomAD v3.1.2, and UK~Biobank.
  Directly applicable to the analysis of genome-wide screen and ’omics data.
  \href{https://github.com/fullscreen-triangle/gospel}{github.com/fullscreen-triangle/gospel}
}

\cvitem{lavoisier}{
  \textbf{HPC Mass-Spectrometry / Metabolomics Pipeline.}
  Dual-pipeline MS analysis (numerical MS1/MS2 + CNN classification, PyTorch);
  memory-mapped I/O for datasets of arbitrary size; Ray distributed execution;
  98.9\,\% NIST accuracy. Used for $>$100\,GB mzML datasets in lipidomics and
  metabolomics --- the data class of the project's metabolomics arm.
  \href{https://github.com/fullscreen-triangle/lavoisier}{github.com/fullscreen-triangle/lavoisier}
}

\cvitem{nebuchadnezzar}{
  \textbf{Computational Pharmacology \& Drug-Metabolism Modelling.}
  Derivation of pharmacokinetic / pharmacodynamic and drug--target trajectories
  from molecular structure; validated 39/39 against NIST reference compounds.
  Relevant to the mechanism and processing of antimetabolite-class compounds.
  \href{https://github.com/fullscreen-triangle/nebuchadnezzar}{github.com/fullscreen-triangle/nebuchadnezzar}
}

%----------------------------------------------------------------------------------------
%	TECHNICAL SKILLS
%----------------------------------------------------------------------------------------

\section{Technical Skills}

\cvitem{Laboratory}{
  Mammalian cell culture; fluorescence microscopy and imaging; molecular and
  cell biology; DNA-damage / repair and crosslinking assays; nucleic-acid and
  protein methods
}

\cvitem{Languages}{
  Python (primary), R, Rust, Bash, SQL, TypeScript
}

\cvitem{Bioinformatics}{
  Variant calling, RNA-seq quantification, CRISPR-screen and pathway analysis
  (KEGG, Reactome, GSEA); MS / metabolomic data analysis (LIPID MAPS, HMDB,
  NIST); FASTQ/BAM/VCF data types; Nextflow pipeline patterns
}

\cvitem{ML / Statistics}{
  PyTorch; scikit-learn, XGBoost; Bayesian-network and mixture modelling;
  high-dimensional statistical analysis
}

\cvitem{Reference data}{
  1000~Genomes, gnomAD, UK~Biobank, IEDB, NIST, LIPID MAPS, HMDB
}

%----------------------------------------------------------------------------------------
%	LANGUAGES
%----------------------------------------------------------------------------------------

\section{Languages}

\cvitemwithcomment{English}{Native / Fluent}{working language of the position}
\cvitemwithcomment{German}{Conversational}{living in Germany since 2011}

%----------------------------------------------------------------------------------------

\renewcommand{\listitemsymbol}{-~}

\end{document}
